% =============================================================
% File: chapters/ta/chapter-3/kebutuhan-sistem.tex
%
% PERUBAHAN:
%   - "sistem audit log" → "mekanisme audit log" di pembuka
%   - Caption tabel KF, KNF, pemetaan: "Audit Log" → "Audit Log"
%   - KNF-TA-03: "data audit log" → "data audit log"
%   - KF-TA-04: "rekaman audit log" → "audit record"
%   - Paragraf satu baris
% =============================================================

\subsection{Kebutuhan Sistem}

Berdasarkan masalah yang telah dipetakan pada Tabel~\ref{tab:identifikasi-masalah}, kebutuhan sistem pada Tugas Akhir ini dirumuskan khusus untuk mengembangkan mekanisme audit log terintegrasi yang mampu mengelola, memvalidasi, dan menyimpan bukti digital dari berbagai komponen DecMed. Kebutuhan ini dirancang untuk menyelesaikan masalah belum terintegrasikannya log aktivitas, keterbatasan cakupan \textit{event}, kerentanan modifikasi log, ketidakstandaran format, serta ancaman \textit{data pruning} pada jaringan DLT.

Kebutuhan sistem dibagi menjadi kebutuhan fungsional dan kebutuhan nonfungsional. Kebutuhan fungsional menjelaskan kemampuan otomatis yang harus disediakan oleh sistem dalam menangkap, memformat, dan menyimpan \textit{audit record}. Sementara itu, kebutuhan nonfungsional menjelaskan batasan kualitas, integritas, dan ketersediaan data log dalam jangka panjang.

\subsubsection{Kebutuhan Fungsional}

Kebutuhan fungsional dirumuskan dari sudut pandang interaksi antar-proses (\textit{system/process actors}) antara komponen internal DecMed dan repositori audit log. Rangkuman kebutuhan fungsional ditunjukkan pada Tabel~\ref{tab:kebutuhan-fungsional}.

\begin{table}[htbp]
    \centering
    \caption{Kebutuhan Fungsional Sistem Audit Log}
    \label{tab:kebutuhan-fungsional}
    \renewcommand{\arraystretch}{1.3}
    \begin{tabular}{|p{2cm}|p{11.5cm}|}
        \hline
        \textbf{ID} & \textbf{Deskripsi Kebutuhan} \\
        \hline
        KF-TA-01 & Sistem harus dapat mengagregasi dan mengelola data log dari berbagai komponen DecMed ke dalam satu repositori audit log terpusat. \\
        \hline
        KF-TA-02 & Sistem harus dapat menangkap dan merekam keseluruhan jenis \textit{event} yang sudah ditentukan pada Tabel~\ref{tab:threat-event-mapping}. \\
        \hline
        KF-TA-03 & Sistem harus memformat dan menstandardisasi \textit{audit record} ke dalam skema data terstruktur sebelum disimpan. \\
        \hline
        KF-TA-04 & Sistem harus menyimpan \textit{audit record} ke dalam media penyimpanan yang bersifat \textit{append-only} untuk mencegah modifikasi (\textit{overwrite}). \\
        \hline
        KF-TA-05 & Sistem menyediakan mekanisme retensi dan pengambilan kembali arsip audit log melalui IPFS pinning dan metadata on-chain \\
        \hline
    \end{tabular}
\end{table}

\subsubsection{Kebutuhan Nonfungsional}

Kebutuhan nonfungsional difokuskan untuk memastikan atribut kualitas keamanan, integritas, dan keandalan data log. Rangkuman kebutuhan nonfungsional ditunjukkan pada Tabel~\ref{tab:kebutuhan-nonfungsional}.

\begin{table}[htbp]
    \centering
    \caption{Kebutuhan Nonfungsional Sistem Audit Log}
    \label{tab:kebutuhan-nonfungsional}
    \renewcommand{\arraystretch}{1.3}
    \begin{tabular}{|p{2cm}|p{11.5cm}|}
        \hline
        \textbf{ID} & \textbf{Deskripsi Kebutuhan} \\
        \hline
        KNF-TA-01 & \textbf{Integritas:} Sistem harus menjamin bahwa data \textit{audit record} yang telah disimpan bersifat \textit{tamper-evident}. \\
        \hline
        KNF-TA-02 & \textbf{Interoperabilitas \& Keterbacaan:} Format data log yang dihasilkan harus konsisten, terstruktur, dan mudah dibaca/diproses oleh komponen sistem. \\
        \hline
        KNF-TA-03 & \textbf{Ketersediaan Data:} Sistem harus menjamin ketersediaan dan keteraksesan data audit log dalam jangka panjang. \\
        \hline
        KNF-TA-04 & \textbf{Konsistensi Pencatatan:} Sistem harus memastikan setiap \textit{event} yang terjadi dicatat secara presisi tanpa ada kejadian bernilai audit tinggi yang terlewat. \\
        \hline
    \end{tabular}
\end{table}

\subsubsection{Pemetaan Masalah terhadap Kebutuhan Sistem}

Untuk memastikan bahwa setiap masalah yang diidentifikasi pada Tabel~\ref{tab:identifikasi-masalah} memiliki solusi yang sesuai, dilakukan pemetaan terhadap kebutuhan fungsional dan nonfungsional. Hasil pemetaan ditunjukkan pada Tabel~\ref{tab:pemetaan-masalah}.

\begin{table}[htbp]
    \centering
    \caption{Pemetaan Masalah terhadap Kebutuhan Sistem Audit Log}
    \label{tab:pemetaan-masalah}
    \renewcommand{\arraystretch}{1.3}
    \begin{tabular}{|p{2.5cm}|p{5cm}|p{5.5cm}|}
        \hline
        \textbf{ID Masalah} & \textbf{Kebutuhan Fungsional} & \textbf{Kebutuhan Nonfungsional} \\
        \hline
        M-TA-01 & KF-TA-01 & KNF-TA-02, KNF-TA-04 \\
        \hline
        M-TA-02 & KF-TA-02 & KNF-TA-04 \\
        \hline
        M-TA-03 & KF-TA-04 & KNF-TA-01 \\
        \hline
        M-TA-04 & KF-TA-03 & KNF-TA-02 \\
        \hline
        M-TA-05 & KF-TA-05 & KNF-TA-03 \\
        \hline
    \end{tabular}
\end{table}

Tabel~\ref{tab:pemetaan-masalah} menunjukkan bahwa seluruh masalah (M-TA-01 hingga M-TA-05) telah ditangani secara langsung oleh kombinasi kebutuhan fungsional dan nonfungsional. Pemetaan ini menjadi dasar penentuan rancangan arsitektur dan skenario pengujian pada bab selanjutnya.